\section{附录 A：定理 \ref{thm:three-value} 的证明细化：返回塔分解与三值上界}
\label{app:proof-thm31}

本附录给出定理 \ref{thm:three-value} 的一条可复写证明链。

\subsection*{A.1 返回域的区间性}

对任意 $n\ge 1$，定义
\[
E_n=\{x\in E:\ r_E(x)=n\}.
\]
由于 $T_\theta$ 为等距同胚且 $E$ 为区间，可证明每个 $E_n$ 为至多一个区间。证明思路为：若 $x<y$ 且二者的首返时间均为 $n$，则对任意 $z\in(x,y)$ 由单调性与等距性可推出 $T_\theta^k z\notin E$（$1\le k<n$）且 $T_\theta^n z\in E$，从而 $z\in E_n$。

因此 $r_E$ 为分段常值函数。

\subsection*{A.2 可能返回时间的限制}

取 $i\ge 0$ 使 $\delta_i<b\le \delta_{i-1}$。该选择意味着 $q_i\theta$ 的模 $1$ 偏移量 $\delta_i$ 小于窗口长度，而上一层偏移 $\delta_{i-1}$ 不小于窗口长度。由连分数逼近的最佳性，可证明：首返时间仅可能在 $\{q_i, q_{i+1}-k q_i\}$ 的有限集合内取值，且 $k$ 的范围由 $a_{i+1}$ 控制。

\subsection*{A.3 参数 $K$ 与三值结构}

定义
\[
K=\max\{k\ge 0:\ k\delta_i+\delta_{i+1}<b\}.
\]
比较 $b$ 与 $k\delta_i+\delta_{i+1}$ 的位置决定“第三塔高”是否出现，以及对应返回域在 $E$ 内的相对位置。由此可将 $E$ 划分为至多三段区间，对应返回时间分别为 $q_i$、$q_{i+1}-(K-1)q_i$、$q_{i+1}-Kq_i$（顺序由奇偶性与端点约定决定）。
同时由 $\delta_{i-1}=a_{i+1}\delta_i+\delta_{i+1}$ 与 $\delta_i<b\le\delta_{i-1}$ 推出 $0\le K\le a_{i+1}-1$。

\subsection*{A.4 加法闭包关系}

由分母递推 $q_{i+1}=a_{i+1}q_i+q_{i-1}$ 与三值表达可得：两侧值与中间值之间满足一次线性关系；在合适排列下化为“中间值等于两侧值之和”。该关系对应返回塔拼接的组合结构。

